\chapter{Probe sức khỏe và quản lý tài nguyên}
\label{ch:probes}

Chương này là trái tim vận hành của Phần~V: kubelet quyết định một pod
còn sống và sẵn sàng như thế nào, và bộ lập lịch cùng kernel chia hai
nhân CPU cho một tá JVM ra sao. Gần như mỗi dòng ở đây đều học được từ
việc nhìn một pod thật cư xử sai.

\section{Ba loại probe}

\begin{description}
  \item[startupProbe] ``vẫn đang boot --- kiên nhẫn.'' Khi nó còn thất
  bại, các probe khác bị tạm hoãn. Cạn ngân sách thì container bị khởi
  động lại.
  \item[readinessProbe] ``có thể nhận traffic \emph{ngay bây giờ}
  không?'' Thất bại thì pod bị rút khỏi endpoints của Service --- không
  restart, không kịch tính, chỉ là không nhận request mới cho tới khi
  qua lại.
  \item[livenessProbe] ``tiến trình đã hết cứu được chưa?'' Thất bại
  thì container bị \emph{khởi động lại}. Cái nguy hiểm.
\end{description}

\begin{figure}[h]
\centering
\begin{tikzpicture}[node distance=4mm and 12mm]
  \node[box] (boot) {khởi động\\\scriptsize startup probe lỗi};
  \node[boxfill, right=of boot] (ready) {Ready\\\scriptsize có trong endpoints};
  \node[box, right=of ready] (notready) {chưa Ready\\\scriptsize bị loại khỏi endpoints};
  \draw[flow] (boot) -- node[lbl,above]{startup đạt} (ready);
  \draw[flow, bend left=18] (ready) to node[lbl,above]{readiness lỗi} (notready);
  \draw[flow, bend left=18] (notready) to node[lbl,below]{readiness đạt} (ready);
  \draw[flow, bend right=25] (boot.south) to node[lbl,below]{hết ngân sách: restart} ++(0,-0.75);
\end{tikzpicture}
\caption{Vòng đời của pod dưới con mắt kubelet.}
\end{figure}

Từ \code{course-service.yaml}:

\begin{lstlisting}[language=yaml]
startupProbe:
  httpGet: { path: /v3/api-docs, port: http }   (*@\co{1}@*)
  periodSeconds: 10
  failureThreshold: 12                          (*@\co{2}@*)
readinessProbe:
  httpGet: { path: /v3/api-docs, port: http }
  periodSeconds: 5
  failureThreshold: 3
# no livenessProbe                              (*@\co{3}@*)
\end{lstlisting}

\begin{description}
  \item[\co{1}] Quyết định về đường dẫn probe. \code{/actuator/health}
  là endpoint sinh ra cho việc này --- nhưng course-service và
  auth-service không cho nó qua Spring Security, nên nó trả 401 và
  kubelet tính là thất bại. Đường dẫn không cần xác thực duy nhất chỉ
  trả 200 \emph{sau khi toàn bộ context đã lên} --- Flyway đã migrate,
  bean đã dựng --- là \code{/v3/api-docs} của springdoc.
  Dictionary-service và gateway \emph{có} cho phép actuator và probe nó
  đúng cách; còn notification-service không có Spring Security nào cả,
  nên endpoint health của nó mặc định mở --- ba service, ba câu trả lời
  cho cùng một câu hỏi. Quy tắc:
  đường dẫn probe phải không cần xác thực \emph{và} là tín hiệu
  readiness thật; một kiểm tra \code{tcpSocket} qua ngay khi Tomcat
  bind cổng thì chẳng đạt điều nào.
  \item[\co{2}] $12 \times 10$\,s = hai phút ân hạn khởi động --- đo
  cho một JVM khởi động nguội trên hai nhân bị tranh chấp.
  \item[\co{3}] Cố tình vắng mặt. Liveness giết; thứ duy nhất nó bắt
  được thêm so với readiness là một JVM treo cứng thật sự --- hiếm ---
  trong khi một đợt dừng GC dài trên máy 2 nhân đang tải nặng có thể làm
  probe thất bại và giết một pod \emph{khỏe mạnh}. Hãy thêm liveness khi
  bạn đã quan sát được cái treo mà nó sinh ra để bắt.
\end{description}

\begin{pitfall}[Startup probe failed: statuscode 404 --- hoặc 403]
Triệu chứng (milestone~6, build \#6): log của notification-service cho
thấy Kafka đã gán partition và ứng dụng khởi động hoàn toàn, vậy mà pod
không bao giờ Ready; probe báo \code{404}. dictionary-service cũng vậy,
nhưng \code{403}. Nguyên nhân: các file config-server \emph{expose}
\code{health,info,metrics,prometheus} cho mọi service --- nhưng chỉ
config-server, discovery-server và api-gateway có
\code{spring-boot-starter-actuator} trên classpath. Cấu hình expose cho
một thư viện không tồn tại thì lặng lẽ vô hiệu; URL đơn giản là 404.
Biến thể 403 là đóng góp của Spring Security: 404 forward sang
\code{/error}, \code{/error} không nằm trong danh sách permitAll, và
một chuỗi filter stateless không có entry point thì trả 403 --- nên
triệu chứng chỉ về phân quyền trong khi nguyên nhân là thiếu jar.
Sửa: thêm starter vào cả bốn POM của service. Bài học: đường dẫn probe
phải được kiểm chứng bằng \code{curl} từ trong pod \emph{trước khi} tin
nó trong manifest --- khẳng định trước đó của sách rằng
dictionary-service ``probe health đúng cách'' được viết từ
SecurityConfig của nó, không phải từ một request thật.
\end{pitfall}

\begin{pitfall}[vòng lặp restart tự sinh ra chính nó]
Triệu chứng: một service khởi động lại vài phút một lần khi tải nặng,
mỗi lần restart làm tải nặng thêm. Nguyên nhân: một liveness probe quá
gắt bị timeout trong lúc GC dừng hoặc gọi downstream chậm --- probe
\emph{gây ra} chính sự cố nó sinh ra để phát hiện. Kiểu lỗi này là lý do
``không có liveness probe'' là mặc định có thể biện hộ, và là lý do
liveness, khi được thêm, chỉ nên kiểm tra ``event loop có còn phản hồi
không,'' không bao giờ kiểm tra phụ thuộc downstream.
\end{pitfall}

\section{Requests, limits và JVM}

\begin{lstlisting}[language=yaml]
resources:
  requests: { memory: 512Mi, cpu: 200m }   (*@\co{1}@*)
  limits: { memory: 768Mi }                (*@\co{2}@*)
env:
  - { name: JAVA_TOOL_OPTIONS, value: "-XX:MaxRAMPercentage=75.0" } (*@\co{3}@*)
\end{lstlisting}

\begin{description}
  \item[\co{1}] \emph{Requests} là sổ sách của bộ lập lịch: pod chỉ
  được đặt vào nơi còn 512\,Mi và 0,2 CPU chưa ai nhận. Requests
  \emph{giữ chỗ}, không giới hạn.
  \item[\co{2}] \emph{Limits} là trần do kernel thực thi. Vượt limit bộ
  nhớ là tiến trình bị OOM-kill --- trạng thái \code{OOMKilled}, exit
  code 137, không Java exception, không dòng log. Để ý sự bất đối xứng:
  có limit bộ nhớ, nhưng \emph{không có limit CPU} --- xem hộp quyết
  định tiếp theo.
  \item[\co{3}] JVM nhận biết container định cỡ heap theo tỷ lệ của
  limit bộ nhớ \emph{của container}: 75\% của 768\,Mi cho heap, phần
  còn lại cho metaspace, thread và off-heap. Không có dòng này, giá trị
  mặc định để lại hoặc lãng phí, hoặc biên an toàn OOM-kill bằng không.
\end{description}

\begin{decision}[vì sao không có limit CPU]
Limit CPU bóp nghẹt (throttle): container vượt hạn mức bị tạm dừng tới
chu kỳ lập lịch tiếp theo --- và JVM đang khởi động chính là lúc service
muốn từng chu kỳ CPU (nạp class, JIT). Khởi động bị bóp biến 60 giây
thành 200 và làm startup probe thất bại: chính limit chế tạo ra thất
bại. Khác bộ nhớ, CPU là tài nguyên nén được --- tranh chấp làm mọi thứ
chậm lại thay vì giết --- nên mẫu dùng ở đây là lời khuyên chuẩn:
\textbf{luôn request CPU, hiếm khi limit}. Bộ nhớ có limit vì kiểu lỗi
của bộ nhớ không giới hạn là kernel giết \emph{một ai đó}, và bạn muốn
ranh giới ấy đoán trước được.
\end{decision}

\begin{hardware}
Tổng requests của mọi service xấp xỉ phần còn lại của 19\,GB sau khi
trừ Jenkins, k3s và các datastore --- một cách có chủ đích. Requests là
lời hứa với bộ lập lịch; hứa quá tay trên một node duy nhất nghĩa là pod
không bao giờ được lập lịch (\code{Pending} mãi mãi), hứa thiếu nghĩa là
chơi roulette OOM. Trên một node, bạn đang hoạch định dung lượng bằng
tay; autoscaler trên cloud biến cùng phép tính đó thành định cỡ node
pool.
\end{hardware}

\begin{handson}[xem một lần khởi động qua con mắt của probe]
\begin{lstlisting}[language=bash]
kubectl -n ts delete pod -l app=course-service
kubectl -n ts get pods -w
# Running 0/1  <- startup probe still failing; no traffic yet
# Running 1/1  <- /v3/api-docs answered; endpoints updated
kubectl -n ts describe pod -l app=course-service | tail -n 12
\end{lstlisting}
Khoảng cách giữa \code{0/1} và \code{1/1} là Flyway cộng khởi động
context --- cửa sổ mà, nếu không có probe, request sẽ đập vào một
service mới lên được một nửa.
\end{handson}

\begin{quiz}
\begin{enumerate}
  \item Với từng probe: thất bại thì chuyện gì xảy ra, và probe nào an
  toàn để đặt gắt?
  \item Vì sao \code{/v3/api-docs} là tín hiệu readiness trung thực cho
  course-service còn \code{tcpSocket: 8082} là lời nói dối? Điều gì
  khiến \code{/actuator/health} không dùng được ở đó, mà dùng được trên
  dictionary-service?
  \item Một pod hiện \code{OOMKilled}, exit 137, log sạch. Giải thích
  cơ chế và vì sao không có Java exception nào. Hai thứ đầu tiên cần
  kiểm tra?
  \item Requests và limits: bộ lập lịch đọc cái nào, kernel thực thi
  cái nào, và vì sao dự án này giới hạn bộ nhớ mà không giới hạn
  CPU?
\end{enumerate}
\end{quiz}
